\documentclass[11pt]{article}
\usepackage[utf8]{inputenc}
\usepackage[T1]{fontenc}
\usepackage{lmodern}
\usepackage{amsmath,amssymb,amsthm}
\usepackage{array,booktabs,longtable}
\usepackage[margin=0.85in]{geometry}
\setlength{\emergencystretch}{3em}

\newtheorem{theorem}{Theorem}
\newtheorem{proposition}[theorem]{Proposition}
\theoremstyle{definition}
\newtheorem{definition}[theorem]{Definition}
\newtheorem{remark}[theorem]{Remark}
\newcommand{\EqSrc}{\ensuremath{\mathrm{EqSrc}}}
\newcommand{\Fun}{\operatorname{Fun}}
\newcommand{\im}{\operatorname{im}}

\title{EQSRC-REPRESENTATIVE-IRRELEVANCE-THEOREM-V1\\
Orbit factorization, invariant source readouts, and scoped countermodels}
\author{Ontology Formalizer packet RT-20260720-024---draft/control}
\date{2026-07-21}

\begin{document}
\maketitle

\begin{abstract}
This packet answers v21 P3-T03 on an explicit source-side target.  For a
nonempty finite representative domain $D_X$, a declared structural action
$G_X\curvearrowright D_X$, and a source-readout family fixed before
evaluation, the joint readout factors through the orbit quotient exactly when
it is constant on every structural orbit.  Equality for every admissible
representative is stronger: the induced quotient readout must itself be
constant.  The maximal invariant real-valued readout algebra and its finite
Reynolds projection are proved explicitly.  Six exhaustively enumerated
controls show that invariant subfamilies factor, while each raw interface
contains a source-side representative-sensitive readout.  The result is a
scoped orbit-factorization theorem plus a scoped raw-interface obstruction,
not physical gauge, operational-observable, ontology-adoption, general
\EqSrc{}, or promotion authority.
\end{abstract}

\section{Authority, question, and fixed source-only inputs}

The artifact is \texttt{draft/control}, \texttt{proposal-only}, and
\texttt{source-extension data}.  It uses the P2-T03 fixed-core relation-image
criterion, the P3-T01 representative-irrelevance evidence contract, and the
P3-T02 finite source-extension controls.  Current ontology does not derive or
adopt the representative domains, structural actions, readout families, orbit
quotients, or any resulting \EqSrc{} law.  Structural automorphisms are not thereby physical gauge transformations.

The exact question is not whether one can call a structural relabeling
``gauge.''  It is whether a readout declared independently of the desired
answer takes equal values on every admissible representative.  The readout
codomain and rule are therefore fixed before enumeration.  Target atlases,
target metrics, benchmark success, candidate reconstructions, validator
status, registry metadata authority, and generated derivatives are excluded
from the premises.

\begin{definition}[Source-readout package]
A bounded source-readout package is
\[
  \mathcal R_X=(D_X,G_X,\mathcal F_X,O_{\mathcal F},q_X),
\]
where $D_X\neq\varnothing$ is a finite set of admissible source-side
representatives, $G_X$ is a finite declared structural automorphism group
acting on $D_X$, $\mathcal F_X$ is a finite family of readouts
$f:D_X\to A_f$ fixed before evaluation,
\[
  O_{\mathcal F}(d)=(f(d))_{f\in\mathcal F_X},
\]
and $q_X:D_X\to D_X/G_X$ is the structural orbit projection.  These are
task-local source readouts.  Calling them operational physical observables
would require a separate source-to-measurement derivation.
\end{definition}

\section{Exact factorization and equality theorems}

\begin{theorem}[Orbit-factorization criterion]
\label{thm:factor}
There is a unique map
\[
  \overline O_{\mathcal F}:D_X/G_X\longrightarrow
  \prod_{f\in\mathcal F_X}A_f
  \quad\text{with}\quad
  O_{\mathcal F}=\overline O_{\mathcal F}\circ q_X
\]
if and only if
\[
  O_{\mathcal F}(g\!\cdot\! d)=O_{\mathcal F}(d)
  \quad\text{for every }g\in G_X,\ d\in D_X.
\]
Equivalently, every member of $\mathcal F_X$ is constant on each structural
orbit.
\end{theorem}

\begin{proof}
If the factorization exists, then $q_X(g\cdot d)=q_X(d)$, so the displayed
equality follows.  Conversely, orbit constancy makes
$\overline O_{\mathcal F}([d]):=O_{\mathcal F}(d)$ independent of the chosen
representative.  Surjectivity of $q_X$ forces uniqueness.
\end{proof}

\begin{theorem}[All-choice equality criterion]
\label{thm:allchoice}
For nonempty $D_X$, every admissible representative has the same joint
readout if and only if $\im O_{\mathcal F}$ is a singleton.  When Theorem
\ref{thm:factor} applies, this is equivalent to
$\overline O_{\mathcal F}$ being constant on all of $D_X/G_X$.
\end{theorem}

\begin{proof}
Equality of all values gives one image point.  A singleton image gives
equality of every pair.  Nonemptiness is required because the image of an
empty domain is empty, not a uniquely realized value.  Under factorization,
surjectivity of $q_X$ identifies the two image sets.
\end{proof}

\begin{remark}[Multiple-orbit guard]
Orbit factorization removes choice only \emph{within} an orbit.  If
$D_X/G_X$ contains two points on which $\overline O_{\mathcal F}$ differs,
then the readout is structurally invariant but not equal for every admissible
choice.  A quotient theorem and a global unique-image theorem are therefore
distinct results.
\end{remark}

\begin{proposition}[Selector-image form]
Suppose an admissible selector domain $\mathcal S$ is nonempty and a declared
interface sends $\sigma\in\mathcal S$ to a source-readout section $K_*(\sigma)$.
All selectors are irrelevant at that interface exactly when
$\lvert\im K_*\rvert=1$.  If $\mathcal S=\varnothing$, the image is empty and
there is no realized unique readout.  On the P2-T03 fixed core this reduces
componentwise to singleton images of the fixed loci under the declared
interface.
\end{proposition}

\begin{proof}
This is Theorem \ref{thm:allchoice} with $D_X=\mathcal S$ and joint map
$O_{\mathcal F}=K_*$.  The final sentence is the registered P2-T03
relation-image criterion on the fixed certified core.
\end{proof}

\section{Equivariance and the invariant real-readout algebra}

\begin{proposition}[Equivariance is not equality]
Let $G_X$ also act on a codomain $A$ and let $O:D_X\to A$ be equivariant.
Then
\[
  O(g\cdot d)=g\cdot O(d),
\]
which does not imply $O(g\cdot d)=O(d)$.  Equality on every orbit holds
exactly when $\im O\subseteq A^{G_X}$.  The codomain-orbit readout
$d\mapsto[O(d)]\in A/G_X$ is invariant automatically and therefore descends
to $D_X/G_X$, but it returns only a structural codomain orbit, not a unique
raw value.
\end{proposition}

\begin{proof}
The first equation is the definition of equivariance.  Comparing it with
orbit constancy shows that equality requires $g\cdot O(d)=O(d)$ for every
$g$ and $d$, equivalently pointwise fixed image.  Quotienting the codomain
makes $[g\cdot O(d)]=[O(d)]$, which proves only the final statement.  A
nontrivial covariant transformation law therefore transports values rather
than erasing representative dependence.
\end{proof}

\begin{theorem}[Maximal invariant real-readout algebra]
\label{thm:reynolds}
For finite $G_X$, the invariant subalgebra
\[
  \Fun(D_X,\mathbb R)^{G_X}
\]
is naturally isomorphic to $\Fun(D_X/G_X,\mathbb R)$.  The operator
\[
  (P_Gf)(d)=\frac{1}{|G_X|}\sum_{g\in G_X}f(g\cdot d)
\]
is an idempotent linear projection from $\Fun(D_X,\mathbb R)$ onto that
invariant subalgebra.  It is the identity on invariant readouts.
\end{theorem}

\begin{proof}
Theorem \ref{thm:factor} gives the algebra isomorphism by pullback along
$q_X$.  For $h\in G_X$, left multiplication permutes the finite summation
index, so $(P_Gf)(h\cdot d)=(P_Gf)(d)$.  If $f$ is invariant, every summand
equals $f(d)$, hence $P_Gf=f$.  Therefore $P_G^2=P_G$ and its image is exactly
the invariant subalgebra.
\end{proof}

The averaging statement is restricted to finite groups and real-valued
readouts.  It is not a license to average labels, topological sectors,
nonconvex states, or physical measurement outcomes without a separately
declared operation and provenance.  Even when $P_Gf$ exists, replacing $f$
by $P_Gf$ changes the readout contract unless that invariant projection was
declared prospectively.  The Reynolds projection is linear and idempotent but
is generally not multiplicative: for the two-point sign readout $s$ under the
$C_2$ swap, $P_Gs=0$ while $P_G(s^2)=1$.  Averaging therefore constructs a
new invariant readout; it does not prove the original readout irrelevant.

There is also an adequacy guard.  Selecting
$\Fun(D_X,\mathbb R)^{G_X}$ and then observing that its members factor is
mathematically correct but partly definitional.  It does not show that this
subalgebra is operationally complete, empirically adequate, or sufficient for
downstream physics.  Those are separate source-to-readout and measurement
burdens.  Likewise, a P3-T02 marked singleton is added source structure on a
changed domain, not representative irrelevance on the original unmarked
domain.

All theorems above are objectwise.  A natural quotient-readout family across
a source category additionally requires the orbit relation to be a congruence
for every admitted source morphism and the descended readouts to commute with
those morphisms.  The finite controls do not establish that family-level
naturality, nor do they discharge the separately declared descent or
equalizer obligations for noninvertible arrows.

\section{Complete finite source-side controls}

The companion JSON record lists every group element, every representative,
and every readout value.  The task-local validator checks closure, inverses,
orbits, every action--value pair, and the stored summary.  The invariant
subfamily is fixed separately from the raw interface; no sensitive readout is
silently discarded after evaluation.

\begin{longtable}{>{\raggedright\arraybackslash}p{0.22\linewidth}
                  >{\raggedright\arraybackslash}p{0.20\linewidth}
                  >{\raggedright\arraybackslash}p{0.24\linewidth}
                  >{\raggedright\arraybackslash}p{0.24\linewidth}}
\toprule
Control & Structural orbits & Factoring source readouts & Sensitive source readout\\
\midrule
\endhead
Orientation torsor & one orbit on $\{-1,+1\}$
& unsigned norm & orientation sign\\
Four-cycle root & one orbit on four vertices
& cycle length; degree multiset & root label; distance to fixed $v_0$\\
Three-line control & one orbit on three lines
& line dimension; other-line count & selected-line label\\
Equal-block partition & one orbit on two block orderings
& unordered block-size multiset & first-block signature\\
Graded cyclic phase & one orbit on four origins
& cyclic grade multiset; cycle length & absolute phase origin\\
Two-component representative & two two-point orbits
& component class; component size & within-component label\\
\bottomrule
\end{longtable}

All six invariant subfamilies factor through their orbit quotients.  In the
first five transitive controls, the invariant joint readout is also constant
on every admissible choice.  In the sixth control, component class separates
the two quotient points, so factorization does not imply all-choice equality.
Every raw interface contains at least one readout that varies within an orbit;
therefore every raw interface is an exact bounded countermodel to full
representative irrelevance.

\section{Result classification and limitations}

The strongest result supported by the packet is
\texttt{scoped\_orbit\_factorization\_theorem\_with\_raw\_interface\_countermodels}.
It establishes a nontrivial quotient-level source-readout subalgebra and a
scoped raw-interface obstruction.  It does not assert that every later
operational observable must include the sensitive readouts, or that no
conservative source extension can supply a lawful quotient target.  Conversely,
the existence of an invariant subalgebra does not prove that it is physically
complete, operationally measurable, dynamically selected, or robust under the
full source-variation class.

The status is \texttt{blocked\_adoption\_open\_continuation}: the current
ontology does not derive or adopt the readout package or a general source
equivalence, while P3-T04 remains a distinct in-scope probability or dynamics
packet.  The exact P3-T03 theorem route is complete and should not be repeated
without a new source-readout family, variation class, theorem, or countermodel.

\section{Distance-to-GR matrix}

\begin{longtable}{>{\raggedright\arraybackslash}p{0.29\linewidth}
                  >{\raggedright\arraybackslash}p{0.61\linewidth}}
\toprule
Burden & P3-T03 status\\
\midrule
\endhead
Source ontology primitives & unchanged; no canonical primitive or readout law adopted\\
Source equivalence \EqSrc{} & scoped orbit-factorization theorem and raw-interface countermodels; general \EqSrc{} remains open\\
RetainH and GenH & unchanged and blocked by missing primitives\\
ObsLoc$_{lc}$ and Resp$_{lc}$ & unchanged; no localization or response law follows\\
$M_{\rm src}$ and $g_{\rm eff}$ & not started by this packet\\
Matter coupling and Einstein equations & not started by this packet\\
Finite-variation robustness & finite exact action enumeration only; general variation robustness remains open\\
Benchmark promotion and Gate Chair status & unchanged and human-gated\\
Current route freeze or hard-fail status & exact P3-T03 raw-interface repetition frozen; P3-T04 remains materially distinct\\
\bottomrule
\end{longtable}

The new mathematical payload consists of the orbit-factorization theorem, the
all-choice equality and multiple-orbit guards, the finite invariant-algebra
projection, and six exhaustive source-side countermodels.  There is no
scientific Distance-to-GR ledger delta.

\section{Conclusion}

A lawful representative-irrelevance claim requires a source-readout family
fixed before evaluation and a proof that its joint map is constant on the
declared representative domain.  Structural orbit invariance is exactly the
condition for quotient factorization, but it is weaker than global equality
when multiple orbits remain.  Equivariance is weaker still unless the image is
pointwise fixed.  The bounded controls exhibit useful invariant source
subfamilies and simultaneous raw representative sensitivity.  Accordingly,
the packet completes P3-T03 with a scoped theorem and scoped obstruction only;
it establishes no physical gauge, operational observable, ontology adoption,
general \EqSrc{}, downstream GR derivation, promotion, or publication result.

\end{document}
